
ment pour disposer de formules du type (1.1), (1.2), et la même démonstration essentiellement devrait s'appliquer sans hypothèse de caractéristique, une fois résolue dans un sens affirmatif la question soulevée au n\textsuperscript{o} 1.

\textbf{3.2.} Lorsque $X$ est lisse sur un corps, on dispose aussi de l'anneau de Chow
\[
 A(X),
\]
et d'une théorie des classes de Chern sur $X$ à valeurs dans $A(X)$ [9], plus fine que la théorie à valeurs dans $\Gtop^\bullet(X)$ qu'on vient d'introduire, -ces dernières classes de Chern étant images des classes de Chern style Chow par l'homomorphisme canonique surjectif (cf. 4.1)
\[
 \phi:A(X)\longrightarrow \Gtop^\bullet(X).
\]
Cet homomorphisme est un isomorphisme mod. torsion, mais pas un isomorphisme en général (4.7). Ceci dit, si $X$ et $Y$ sont lisses, l'inclusion $f:Y\to X$ définit encore un homomorphisme de degré $d$ de groupes gradués
\[
 f_*:A(Y)\longrightarrow A(X),
\]
ce qui permet de donner un sens à la formule (3.1) comme une formule dans $A^i(X)$\SGAAuthorFootnote{*}{La formule, ainsi que la variante signalée dans 3.3 a été prouvée depu par P. JOUANOLOU (nov.\ 1969). Cf.}. Cette formule se réduit d'ailleurs à $0=0$ pour $i<d$, et ne présente donc d'intérêt que pour $i\geq d$. Même pour $i=d=2$ cette formule n'est pas prouvée à l'heure actuelle (cf. 4.3), ni la formule moins précise qui s'en déduit en remplaçant l'équivalence rationelle des cycles algébriques par l'équivalence algébrique. Contrairement à la formule analogue dans $\Gtop^\bullet(X)$, il ne semble pas y avoir de raison bien convainquante pour supposer que la formule sans dénominateurs doive être vraie dans $A(X)$. Une bonne façon de la tester, vue notre ignorance quasi totale sur la nature des groupes $A^i(X)$ pour $i\geq 2$, semble être de tester les formules qu'on peut en déduire dans des $A^1(Z)=\Pic(Z)$, par exemple en transformant les deux membres de (3.1) par $g_*:A(X)\longrightarrow A(Z)$, où $g:X\to Z$ est un morphisme propre dans un $Z$ quasi-projectif lisse de dimension $\dim X-i+1$.

\textbf{3.3.} Enfin, pour donner un sens à (3.1), on peut également utiliser la théorie cohomologique des classes de Chern (SGA 5 VII), chaque fois qu'on dispose d'\uline{homomorphismes de Gysin}
\[
 f_*:H^p(Y,\mu_n^{\otimes i})\longrightarrow H^{p+2d}(X,\mu_n^{\otimes(i+d)})
\]
(cf. n\textsuperscript{o}~6), ce qui est le cas tout au moins si $Y$ et $X$ sont tous deux lisses sur un schéma $S$ (grâce au théorème de pureté cohomologique relatif 
